\textbf{Versão parcial: só a variante (iii)} (desacoplado, esqueleto rígido, ar passivo: as hipóteses do artigo); as
variantes (ii) e (i) estão rodando e entram na próxima versão. Os três solos de Rahardjo et al.\ (2007) (f10, f50 e
f100) foram calculados com três chuvas de 24~h cada e o FS (SRM) de hora em hora até 48~h. Com as hipóteses do
artigo, o FS mínimo fica a no máximo 0,02 da Eq.~6 do artigo nos nove casos, e no f50 a curva FS $\times$ $t$ passa
sobre os pontos da Fig.~5b (diferença $\le 0{,}03$, a incerteza da digitalização) em 6, 12, 24 e 48~h. O relatório de
outubro ficava 0,03 a 0,17 acima nos mesmos pontos. Rompem (FS $< 1$) o f10 com 180~mm/h (de 11 a 26~h, mínimo 0,842)
e o f50 com 90~mm/h (de 20 a 32~h, mínimo 0,832). O FS mínimo cai no fim da chuva em sete casos; nos dois do f100 com
chuva fraca há um patamar depois da chuva, com o mínimo 0,008 abaixo do valor de 24~h, dentro da resolução do SRM. A
concordância com o artigo ainda não é validação. O FS do SRM ainda não convergiu com o refinamento adaptativo (o
último nível muda 0,003 a 0,012; o convergido deve ficar $\approx 0{,}01$ abaixo), a malha do transiente é grossa
(catetos de 2,5~m; a frente de umedecimento do f100 tem 1 a 2~m) e a chuva entra por metro inclinado (41\,\% mais
água na face que na projeção horizontal); os três pontos são testados na campanha do fermi. O custo é quase todo do
fator: o transiente de 48~h leva $\approx 2$~min e as 49 análises do SRM, 30 a 50~min por caso.
